\chapter{Results}
\label{ch:results}

\begin{quote}\small\itshape
A note on status and provenance. All result numbers are centralised as macros in
\file{results/numbers.tex}; the chapter sources cite macros only. The two-year single-agent
run \runId\ remains the primary reasoning artifact. The like-for-like performance comparison
is now the \commonRunLabel: \commonStart\ to \commonEnd, \commonUniverse, symmetric
$10$\,bps one-way costs, cache-locked prices, timestamped news enabled, and the same decision
window for the single agent and the multi-agent Portfolio Manager.
\end{quote}

\section{Experimental setup}
\label{sec:setup}

The canonical single-agent run is \runId, over \studyStart\ to \studyEnd\ on the universe
\universeList, producing \nDecisions\ decisions with \nParseErrors\ parse errors. It is used
for the primary regime-segmented reasoning diagnosis because it spans the broader period. The
clean financial comparison is separate and explicitly labelled: \commonRunLabel,
\commonStart\ to \commonEnd, \commonUniverse, \commonNTradingDays\ trading days,
\commonSingleDecisions\ single-agent decisions and \commonPmDecisions\ PM decisions.

This common-reference run resolves the two comparability problems found during drafting. First,
the single agent, PM, and baselines all pay the same $10$\,bps one-way transaction cost and all
report turnover and cost drag. Second, the PM is no longer only a one-month validation: it is
run over the same one-year decision window and ticker set as the single agent. Prices are served
from a prewarmed cache covering the decision window plus the 420-calendar-day regime warm-up;
both paths have zero warm-up exclusions. Confidence intervals use an i.i.d.\ bootstrap of daily
returns with \bootstrapResamples\ resamples (\cref{sec:bootstrap}). PM decisions with
\code{is\_error} are excluded from reasoning tables and counted separately.

\section{Financial performance}
\label{sec:resultsfin}

\Cref{tab:commonfinancial} is the unified comparison table. It is the only table in this thesis
that should be read as a direct single-agent versus PM versus baseline comparison.

\begin{table}[t]
  \centering
  \caption[One-year common-reference comparison]{\commonRunLabel\ over \commonStart\ to
  \commonEnd\ on \commonUniverse. All rows use the same decision window, the same tickers, and
  symmetric $10$\,bps one-way costs. PM performance is reported but not headlined because its
  Sharpe confidence interval crosses zero.}
  \label{tab:commonfinancial}
  \scriptsize
  \resizebox{\textwidth}{!}{%
  \begin{tabular}{l S S c S S S}
    \toprule
    Strategy & {Ann. ret. (\%)} & {Sharpe} & {Sharpe $95\%$ CI}
             & {MaxDD (\%)} & {Turnover (\%)} & {Cost drag (bps)} \\
    \midrule
    Single agent & \commonSingleAnnReturn & \commonSingleSharpe & \commonSingleSharpeCI
      & \commonSingleMaxDD & \commonSingleTurnover & \commonSingleCostDrag \\
    PM multi-agent & \commonPmAnnReturn & \commonPmSharpe & \commonPmSharpeCI
      & \commonPmMaxDD & \commonPmTurnover & \commonPmCostDrag \\
    Eq.-wt B\&H & \commonBhAnnReturn & \commonBhSharpe & \commonBhSharpeCI
      & \commonBhMaxDD & \commonBhTurnover & \commonBhCostDrag \\
    Momentum & \commonMomAnnReturn & \commonMomSharpe & \commonMomSharpeCI
      & \commonMomMaxDD & \commonMomTurnover & \commonMomCostDrag \\
    Mean-rev & \commonMrAnnReturn & \commonMrSharpe & \commonMrSharpeCI
      & \commonMrMaxDD & \commonMrTurnover & \commonMrCostDrag \\
    \bottomrule
  \end{tabular}%
  }
\end{table}

The single agent produces a positive, statistically non-zero Sharpe on this costed one-year
window (\commonSingleSharpeCI), but it does not beat equal-weight Buy-\&-Hold on return or
Sharpe. The PM's point return is higher than the single agent's, but its drawdown and turnover
are much larger and the Sharpe interval \commonPmSharpeCI\ crosses zero; it is therefore
reported as an architecture extension, not as a track record. This is the intended reading:
the primary thesis claim is not superior trading performance, but a clean harness that exposes
reasoning, grounding, costs, and temporal leakage on the same artifact.

\Cref{tab:financial} reports the older two-year single-agent context against the three
baselines. It is retained for the reasoning narrative, not as the clean cost-symmetric
comparison.

\begin{table}[t]
  \centering
  \caption[Two-year single-agent context]{Historical two-year single-agent context over
  \studyStart\ to \studyEnd. This table supports the reasoning diagnosis over a broader
  regime mix. The direct cost-symmetric comparison is \cref{tab:commonfinancial}.}
  \label{tab:financial}
  \begin{tabular}{l S S S S}
    \toprule
    Metric & {Agent} & {Eq.-wt B\&H} & {Momentum} & {Mean-rev} \\
    \midrule
    Annualised return (\%) & \agentAnnReturn & \bhAnnReturn & \momAnnReturn & \mrAnnReturn \\
    Sharpe                 & \agentSharpe   & \bhSharpe    & \momSharpe    & \mrSharpe \\
    Sharpe $95\%$ CI       & \multicolumn{1}{c}{\agentSharpeCI} & {--} & {--} & {--} \\
    Sortino                & \agentSortino  & {--} & {--} & {--} \\
    Calmar                 & \agentCalmar   & {--} & {--} & {--} \\
    Max drawdown (\%)      & \agentMaxDD    & {--} & {--} & {--} \\
    Hit rate (\%)          & \agentHitRate  & {--} & {--} & {--} \\
    \bottomrule
  \end{tabular}
\end{table}

In the older two-year context, the agent nearly matches equal-weight Buy-\&-Hold on raw annualised return
(\agentAnnReturn\,\% versus \bhAnnReturn\,\%) while edging it on the Sharpe ratio
(\agentSharpe\ versus \bhSharpe) at a contained maximum drawdown of \agentMaxDD\,\%, and it
clears the systematic momentum and mean-reversion baselines on both return and Sharpe. The
reading is that the agent earns its keep through \emph{risk control} rather than raw return
generation, consistent with a disciplined, policy-capped sizing strategy --- though the
agent's Sharpe confidence interval \agentSharpeCI\ overlaps the Buy-\&-Hold point estimate, so
the edge is not statistically separated on a single run.

\section{Regime composition of the canonical run}
\label{sec:resultsregime}

Per-regime risk-adjusted decompositions of the equity curve were not computed for the canonical
run (the run records aggregate financial metrics and per-decision regime labels, not per-regime
return series); a regime-resolved equity attribution is noted as follow-up in
\cref{sec:future}. The thesis's regime-resolved analysis is therefore carried by the
\emph{reasoning} metrics (\cref{sec:resultsreasoning}), which are segmented per decision.
\Cref{tab:regimecomp} gives the regime composition that those segmented tables draw on, so the
per-regime sample sizes can be read with appropriate caution.

\begin{table}[t]
  \centering
  \caption[Regime composition of the canonical run]{Regime composition of the canonical
  single-agent run under the causal labels of \cref{sec:regime}. The bull regime dominates the
  two-year window, with substantial bear and high-volatility minorities --- the spread that
  makes regime-segmented reasoning evaluation meaningful. The action mix
  (\actBuy\ buys, \actSell\ sells, \actHold\ holds) foreshadows the central finding: trading is
  overwhelmingly holds, and the rare sells are where the exit-discipline gap appears.}
  \label{tab:regimecomp}
  \begin{tabular}{l S[round-precision=0] S[round-precision=1]}
    \toprule
    Regime & {$n$ decisions} & {\% of run} \\
    \midrule
    Bear           & \nBear    & \pctBear \\
    Bull           & \nBull    & \pctBull \\
    High volatility& \nHighVol & \pctHighVol \\
    Range          & \nRange   & \pctRange \\
    \midrule
    Total          & \nDecisions & {100.0} \\
    \bottomrule
  \end{tabular}
\end{table}

\section{Reasoning quality}
\label{sec:resultsreasoning}

\subsection{Grounding}

\begin{table}[t]
  \centering
  \caption[Grounding by regime]{Grounding of numeric claims, by regime, for the canonical run.
  Grounding is computed deterministically (\cref{sec:grounding}) by matching each numeric claim
  in the rationale against the recorded tool outputs within a $1\%$ tolerance. The near-perfect
  score (\groundingNumGrounded\ of \groundingNumClaims\ claims grounded overall) establishes
  that the agent's \emph{factual} citations are real --- the precondition for reading any
  faithfulness gap as a genuine knowledge--action gap rather than mere hallucination.}
  \label{tab:grounding}
  \begin{tabular}{l S}
    \toprule
    Regime & {Grounding} \\
    \midrule
    Bear            & \groundingBear \\
    Bull            & \groundingBull \\
    High volatility & \groundingHighVol \\
    Range           & \groundingRange \\
    \midrule
    Overall         & \groundingOverall \\
    \bottomrule
  \end{tabular}
\end{table}

Grounding is essentially perfect in every regime (\groundingOverall\ overall), which is the
desirable outcome: the agent does not fabricate the numbers it reasons over. This is what
licenses the interpretation of the faithfulness results that follow.

\subsection{Faithfulness: v1 versus v2 by regime}

\Cref{tab:faithfulness} reports the policy-aware (v2) faithfulness alongside the v1 judge, by
regime. The side-by-side presentation is itself a finding: it shows how much of the apparent
unfaithfulness under the naive judge was an artifact of ignoring the agent's stated sizing
policy (\cref{sec:v2judge}). Across the run, v2 raises faithfulness from \faithVoneOverall\ to
\faithVtwoOverall, and --- because the policy-aware judge resolves capacity cases within its
prompt --- the constrained category falls to zero, so the constrained bucket is no longer doing
hidden work.

\begin{table}[t]
  \centering
  \caption[Faithfulness by regime, v1 versus v2]{Faithfulness by regime under the v1
  (max-deploy) and v2 (policy-aware) judges, canonical run. v2 is the canonical metric; v1 is
  retained for comparison and reproducibility. ``v2 unfaithful'' counts the genuine
  knowledge--action gaps under v2.}
  \label{tab:faithfulness}
  \begin{tabular}{l S[round-precision=0] S S S[round-precision=0]}
    \toprule
    Regime & {$n$} & {v1 faith.} & {v2 faith.} & {v2 unfaithful} \\
    \midrule
    Bear            & \nBear    & \faithVoneBear    & \faithBear    & \nUnfBear \\
    Bull            & \nBull    & \faithVoneBull    & \faithBull    & \nUnfBull \\
    High volatility & \nHighVol & \faithVoneHighVol & \faithHighVol & \nUnfHighVol \\
    Range           & \nRange   & \faithVoneRange   & \faithRange   & \nUnfRange \\
    \midrule
    Overall         & \nDecisions & \faithVoneOverall & \faithVtwoOverall & \nUnfaithfulOverall \\
    \bottomrule
  \end{tabular}
\end{table}

\noindent The qualitative pattern is clear and consequential. Faithfulness is \emph{lowest in
the bull regime} (\faithBull) and highest in range (\faithRange), with bear (\faithBear) and
high volatility (\faithHighVol) in between. The agent's reasoning is least aligned with its
actions precisely when prices are rising: it states caution or a trim condition and keeps the
position. Combined with near-perfect grounding, this is the knowledge--action gap in its
clearest form --- the agent sees and correctly cites the evidence, then acts against it. The
dominant mechanism is examined next.

\subsection{The dominant unfaithfulness driver}

The largest single class of unfaithful decisions is \emph{non-trimming}: cases where the judge
predicts a sell/trim and the agent holds a position that has grown past the cap. We report this
class because the v1 judge silently absorbed it into ``constrained,'' hiding it entirely; the
policy-aware v2 judge exposes it (\cref{sec:v2judge}).

\begin{table}[t]
  \centering
  \caption[Non-trimming as the dominant unfaithfulness class]{Decomposition of the
  \nUnfaithfulOverall\ unfaithful decisions (canonical run). Non-trimming dominates, which
  reframes the agent's weakness specifically as \emph{exit/trim discipline after near-cap
  entries} rather than a diffuse reasoning failure. The concentration figure is the share of
  non-trim cases in the single strongest-trending name.}
  \label{tab:nontrim}
  \begin{tabular}{l S}
    \toprule
    Quantity & {Value} \\
    \midrule
    Non-trim sells (judge=sell, agent=hold) & \nNonTrim \\
    \quad as share of all unfaithful (\%)   & \pctNonTrim \\
    True entry gaps (below target, bullish, cash available) & \nDoneGaps \\
    Concentration in the dominant ticker (\%) & \nvdaNonTrimShare \\
    \bottomrule
  \end{tabular}
\end{table}

\noindent Non-trim sells account for \nNonTrim\ of the \nUnfaithfulOverall\ unfaithful decisions
(about \pctNonTrim\,\%), and roughly \nvdaNonTrimShare\,\% of them are concentrated in the single
strongest-trending name. True entry gaps --- positions held below target with cash available and
a bullish signal --- number only \nDoneGaps. The agent's reasoning failure is thus specific and
interpretable: it builds positions toward the cap and then does not cut them when the evidence
turns, rather than failing diffusely across decision types.

\section{Validity: judge--human agreement}
\label{sec:resultsanchor}

The faithfulness judge is anchored to human annotation (\cref{sec:anchor}). Because the
annotation sample is stratified rather than representative, agreement is reported
\emph{per stratum} in \cref{tab:agreement}; a single global figure would misrepresent a
deliberately rare-case-weighted sample.

\begin{table}[t]
  \centering
  \caption[Judge--human agreement]{Judge--human action agreement on the stratified annotation
  sample ($n=\nAnnotated$). Per-stratum agreement is the meaningful quantity; the global
  figures are reported with the explicit caveat that the sample over-represents diagnostic
  strata and is not a population estimate.}
  \label{tab:agreement}
  \begin{tabular}{l S}
    \toprule
    Quantity & {Value} \\
    \midrule
    Annotated rows & \nAnnotated \\
    Global action agreement (\%) & \judgeHumanAgreement \\
    Global Cohen's $\kappa$ & \judgeHumanKappa \\
    \midrule
    Agreement on non-trim stratum (\%) & \agreementNonTrim \\
    Agreement on entry-gap stratum (\%) & \agreementDone \\
    \bottomrule
  \end{tabular}
\end{table}

\noindent Global agreement is \judgeHumanAgreement\,\% ($\kappa=\judgeHumanKappa$, substantial).
The per-stratum view is more informative than the headline: agreement is \emph{perfect}
(\agreementNonTrim\,\%) on the dominant non-trim class --- the finding that matters most ---
and \emph{poor} (\agreementDone\,\%) on the small entry-gap stratum, where decisions sit around
the $10\%$ target threshold under weak or mixed signals and even a careful human disagrees with
the judge. A separate \emph{blind} reconciliation on the medium window --- annotations filled
against sealed fields, the key opened only afterward --- gave an action agreement of
\blindAgreement\,\% with $\kappa=\blindKappa$ over $n=\blindN$ decisions. The two procedures
together support the v2 judge for the thesis, with the residual weakness localised and
documented rather than hidden.

\section{Multi-agent Portfolio Manager (common-reference extension)}
\label{sec:resultspm}

The PM results below are from the same one-year common-reference run used in
\cref{tab:commonfinancial}: \commonPmRunId, \commonStart\ to \commonEnd, \commonUniverse.
They characterise the multi-agent pipeline's behaviour on the shared benchmark. Because the
Sharpe interval crosses zero, the PM remains an architecture extension rather than a headline
performance claim.

\begin{table}[t]
  \centering
  \caption[Portfolio-Manager evaluation (common-reference run)]{Evaluation of the multi-agent
  Portfolio Manager on the common-reference window. PM faithfulness is measured at the
  allocation level (\cref{sec:pmeval}); the strict variant excludes capacity-constrained
  decisions. Evidence grounding audits analysts' cited evidence against tool outputs. The MCP
  time-machine audit confirms no future timestamp reached the agent.}
  \label{tab:pm}
  \begin{tabular}{l S}
    \toprule
    Quantity & {Value} \\
    \midrule
    PM decisions & \commonPmDecisions \\
    Errored decisions ($n_{\text{pm\_errors}}$) & \commonPmErrors \\
    PM faithfulness (overall) & \commonPmFaith \\
    PM faithfulness (strict) & \commonPmFaithStrict \\
    Analyst evidence grounding & \commonPmGrounding \\
    MCP time-machine violations & \multicolumn{1}{c}{\commonTimeMachineViolations} \\
    Cost drag (bps) & \commonPmCostDrag \\
    Turnover (\%) & \commonPmTurnover \\
    \bottomrule
  \end{tabular}
\end{table}

\noindent Three things stand out, all of them about \emph{correctness} rather than performance.
First, the artifact is leak-free: the MCP time-machine audit found
\commonTimeMachineViolations\ timestamp violations across all checks, with every price read
served from cache and every news item from the corpus. Second, analyst evidence grounding is
effectively perfect (\commonPmGrounding): every checkable technical, risk, and news claim is
traceable to a tool output, and --- following the news-hallucination fix of \cref{sec:rigour} --- news evidence
carries a source and publication date. (The PM's \emph{final} rationale is concise and carries no
numeric claims, so its numeric grounding is reported as ``no claims,'' which is expected, not a
failure.) Third, faithfulness is mixed: strict PM faithfulness is \commonPmFaithStrict, while
overall faithfulness is \commonPmFaith\ once constrained allocations are counted. The three
PM API errors are explicit, excluded from reasoning tables, and included in the reported error
rate; they are system reliability observations, not silent all-cash successes.

We deliberately do \emph{not} headline the PM's annualised return or Sharpe: even on the full
common-reference year, the bootstrap Sharpe confidence interval \commonPmSharpeCI\ crosses
zero. The honest conclusion is that the multi-agent pipeline generalises the evaluation harness
to a richer architecture and remains leak-free and evidence-grounded, but it does not establish
a statistically reliable performance edge over the single agent or Buy-\&-Hold.

\section{News integration: actual state}
\label{sec:resultsnews}

Consistent with \cref{sec:news}, the news corpus is disabled by default and enabled explicitly
for the common-reference runs. The corpus populated for the validation window
(\newsCorpusWindow) is sourced from \newsCorpusSource, with row counts
AAPL${}=\newsCountAAPL$, MSFT${}=\newsCountMSFT$, NVDA${}=\newsCountNVDA$. Provider limitations
are reported honestly: a free Finnhub plan returned no rows for this older window, and a
crawl-based enrichment key failed to authenticate locally, so the corpus relies on the source
that exposes a true publication timestamp. With news enabled, the artifact-level audit found no
\code{published\_at} after \tnow\ and every non-empty news report carried a source-and-date
reference, confirming the evidence discipline of \cref{sec:news} in practice.

\section{Systems and cost}
\label{sec:resultssystems}

On the systems axis, the common-reference single-agent and PM runs consumed
\commonTotalTokens\ tokens and \commonTotalCost\ of incremental model spend, over
\commonWallClockMinutes\ minutes of wall-clock time. The local telemetry recorded
\commonSingleApiCalls\ single-agent/judge API calls and \commonPmApiCalls\ PM-chain API calls,
including cache hits and per-span latency. This local accounting is used because Langfuse OTLP
export is best-effort and timed out on several batches during the long PM run. The test suite
stands at \nTests\ passing tests, the standing evidence for the rigour claim of
\cref{sec:rigour}.

\begin{prelim}
The one-year common-reference run reproduces the qualitative single-agent pattern seen at full
scale --- grounding remains high, faithfulness is strongest in bear and weakest in bull, and
cost/turnover accounting is now explicit. Earlier one-month medium runs are retained only as
development checks and are not thesis performance results.
\end{prelim}
